\section{Methodology}
\label{sec:method}

\subsection{Problem statement and design}
\label{sec:setup}

We test an evaluation claim about a fixed operating point, so the right instrument is an
existing blind benchmark rather than a new prediction run. Let $s_m(r) \in [0,1]$ be method
$m$'s disorder score for residue $r$ and $y(r) \in \{0,1\}$ the experimental label. The rule
under audit is the classifier $\hat{y}(r) = \mathbb{1}[\plddt(r) < 50]$, which under the CAID
encoding $s_{\afplddt} = 1 - \plddt/100$~\citep{piovesan2022} is identically
$\mathbb{1}[s(r) > 0.5]$. Nothing has to be re-predicted: the CAID archives contain both the
ground truth and the official \afplddt submission. We downloaded AlphaFold2 coordinates
separately only to characterise the errors structurally and to verify that the CAID scores
really are $1 - \plddt/100$.

Four design decisions each answer a specific way the result could have been an artefact.
\emph{(i)} Specificity depends entirely on how negatives are defined, so both CAID references
are reported for every result and neither is treated as ``the'' answer. \emph{(ii)} AlphaFold
covers only 82--86\% of CAID2 targets while other methods cover $\approx$100\%, so every
comparison runs on the intersection of targets covered by all panel methods; the \cone/\ctwo
tests are additionally repeated on \plddt's own coverage and on the full reference with missing
targets scored as all-ordered. \emph{(iii)} Comparing all methods at a shared 0.5 cut is
meaningful only for \plddt, so every method is evaluated at three operating points.
\emph{(iv)} Choosing a threshold on the data one then scores is optimistically biased, so the
optimism is measured directly by repeated cross-validation and reported.

\subsection{Data}

We use the four CAID disorder reference sets (\tabref{tab:data}). The two conventions differ
only in the definition of a negative: \dpdb counts as negative only residues observed in an
experimental PDB structure, excluding X-ray-missing residues (hence the large excluded column);
\dnox counts every non-annotated residue as negative. We identified which CAID release each
downloaded reference corresponds to by exact matching on (\#sequences, positives, negatives,
undefined) against the four official reference-set manifests; CAID3 turns out to be the v3
round, so all official comparisons use the v3 files.

\begin{table}[t]
    \centering
    \small
    \begin{tabular}{@{}lrrrrl@{}}
        \toprule
        \textbf{Reference set} & \textbf{Targets} & \textbf{Positives} & \textbf{Negatives} &
        \textbf{Excluded} & \textbf{Release} \\
        \midrule
        \caidtwo \dpdb   & 348 & 37,072 & 93,805 & 156,143 & CAID2 \\
        \caidtwo \dnox   & 210 & 31,315 & 129,487 & 0 & CAID2 \\
        \caidthree \dpdb & 319 & 31,401 & 67,838 & 61,183 & CAID3 v3 \\
        \caidthree \dnox & 204 & 26,367 & 73,610 & 541 & CAID3 v3 \\
        \bottomrule
    \end{tabular}
    \caption{The four evaluation references. Residue counts are as published by CAID. The
    ``Excluded'' column holds residues with no evidence either way; \dpdb discards them, \dnox
    treats almost all of them as ordered. This single convention difference drives the
    specificity result in \secref{sec:c2}.}
    \label{tab:data}
\end{table}

We additionally re-parsed 588 \afdb models (94\% of mappable CAID targets) into per-residue
\plddt and Shrake--Rupley relative solvent accessibility for the structural error analysis.

\subsection{Comparator panel and operating points}

\para{Panel.} We restrict the panel to \emph{general} intrinsic-disorder predictors: 37 methods
on \caidtwo and 56 on \caidthree. It spans classical
predictors (IUPred3~\citep{erdos2021}, ESpritz-D~\citep{walsh2012},
MobiDB-lite~\citep{necci2017,piovesan2025}, VSL2, RONN, DisEMBL), strong sequence-based methods
(flDPnn~\citep{hu2021}, AIUPred, DISOPRED3~\citep{jones2015},
SPOT-Disorder2~\citep{hanson2019}, Metapredict~\citep{emenecker2021}) and protein-language-model
methods (PUNCH2, SETH~\citep{ilzhofer2022}, DisorderUnetLM, ESMDisPred, Dispredict3). Methods that predict a different target --- disordered
\emph{binding} regions (ANCHOR2, MoRFchibi, DisoRDPbind, \dots) or linkers --- are excluded,
because including them would flatter \plddt. Two entries (FoldUnfold, NeProc-disorder) submitted
binary states with an empty score column and cannot be threshold-calibrated; they are excluded
and the exclusion is logged. Excluding them is the conservative choice, since at their implicit
0.5 cut they would have added two artificially weak comparators to the \cfour test. AlphaFold-derived
entries (\afplddt, \afrsa, \afthreeplddt, \afthreersa) are tracked separately as the object of
study and its mechanistic controls, leaving 35 and 52 \emph{dedicated} predictors respectively.

\para{Three operating points.} Every method is scored at (a) the \textbf{literal} threshold
$0.5$, which is the only point at which the folk rule is defined; (b) a \textbf{calibrated}
threshold maximising MCC on the reference, which is our primary setting; and (c) a
\textbf{content-matched} threshold reproducing the true number of disordered residues, following
\citet{kurgan2023}. The literal comparison is reported precisely to show how misleading it is.

\subsection{Metrics and statistics}

We compute residue-level sensitivity, specificity, balanced accuracy (\bacc), MCC, precision,
F1 and ROC-AUC, pooled over residues following CAID's ``dataset'' convention. Uncertainty comes
from bootstrap resampling \emph{over targets} (1000 resamples for binary metrics, 500 for AUC),
which is CAID's own resampling unit and respects the fact that residues within a protein are not
independent. All method-vs-\plddt comparisons use a paired bootstrap with identical target
indices for both methods, so the interval is on the difference, with a two-sided bootstrap
$p$-value. Clause tests \cone and \ctwo are one-sided bootstrap probabilities, \eg
$P(\text{sensitivity} \ge 0.80)$ over resamples.

\para{Two definitions of ``balanced accuracy on short segments''.} The folk claim does not define
this quantity, so we compute two and report both. The \textbf{global-negative} (strict)
definition pairs recall over positive residues in 1--19 aa regions with specificity over
\emph{all} evaluated negatives, so a method cannot buy short-region recall with false positives
elsewhere. The \textbf{locally-flanked} (generous) definition scores each region only against
the ordered residues immediately flanking it (up to $\lceil L/2 \rceil$ evaluated negatives per
side), which removes the influence of a method's global false-positive rate entirely. We also
report region-level detection rates --- a region counts as detected when $\ge 25/50/75\%$ of its
evaluated residues are predicted disordered --- since the claim is phrased about regions.

\subsection{Pipeline validation}
\label{sec:validation}

Before running any experiment we validated the pipeline against the official CAID outputs.
For every full-coverage method our pooled AUC reproduces the published value to three decimals
on all four references (\eg IUPred3 0.8854 vs 0.885, ESpritz-D 0.8989 vs 0.899, MobiDB-lite
0.8680 vs 0.868, SETH-1 0.9115 vs 0.911 on \caidtwo \dpdb; flDPnn 0.8609 vs 0.861 on \caidthree
\dpdb). One trap is worth flagging: the official timing file is keyed by software \emph{group}
and reports the group's best member, so the \textsc{AlphaFold-disorder} group value of
0.947--0.95 is \afrsa (ours 0.9443/0.9498), not \afplddt (ours 0.9287/0.9342). Our threshold
sweep matches the official per-threshold specificity to three decimals at every probed
threshold. Finally, parsing \plddt out of the \afdb coordinates and applying $1-\plddt/100$
reproduces the CAID scores at pooled Pearson $r = 0.9925$ over 323,327 residues in 567 targets
(median per-target $r=0.9955$; median residual 1.14 \plddt units), confirming that the submitted
baseline is the standard \plddt. The residual is \afdb release drift between the CAID submission
and our 2026 snapshot; all benchmark results use the official CAID prediction files, so it does
not propagate.

\subsection{Reproducibility}

Python 3.12.8 with numpy 2.5.3, scipy 1.18.1, pandas 3.0.5, scikit-learn 1.9.0, biopython 1.88.
Seed 42 throughout; all bootstrap index matrices are seeded and shared across methods so that
paired comparisons use identical resamples. No model is trained or run --- all predictions were
downloaded --- so the analysis is CPU-only and completes in under 20 minutes of wall clock.
Re-running the length-stratified and mechanism experiments end to end produced byte-identical
output files.
